\documentclass[11pt]{amsart}
\usepackage[T1]{fontenc}
\usepackage{amsmath,amssymb,amsthm,mathtools,mathrsfs}
\usepackage[margin=1in]{geometry}
\usepackage[colorlinks=true,linkcolor=blue,urlcolor=blue,citecolor=blue]{hyperref}
\emergencystretch=1em
\newtheorem{theorem}{Theorem}[section]
\newtheorem{lemma}[theorem]{Lemma}
\newtheorem{corollary}[theorem]{Corollary}
\theoremstyle{definition}
\newtheorem{definition}[theorem]{Definition}
\newtheorem{remark}[theorem]{Remark}
\newcommand{\Z}{\mathbb Z}
\newcommand{\R}{\mathbb R}
\newcommand{\E}{\mathbb E}
\newcommand{\Var}{\operatorname{Var}}
\newcommand{\TV}{\operatorname{TV}}
\newcommand{\gap}{\operatorname{gap}}
\newcommand{\dist}{\operatorname{dist}_\infty}
\newcommand{\Lip}{\operatorname{Lip}}
\newcommand{\cD}{\mathcal D}
\newcommand{\1}{\mathbf 1}
\newcommand{\Om}{\Omega}
\title{A finite-state proof from window coupling to uniform box gaps}
\author{MOSPROOF appendix}
\date{November 16, 2026}
\begin{document}
\begin{abstract}
For a finite alphabet and a real, translation-invariant, range-one interaction
on $\Z^2$, the window estimate of Definition 4.7 of the manuscript implies a
positive heat-bath spectral gap, uniformly in the boundary condition and in
all sufficiently large rectangular boxes. We give a direct proof, replacing
both cited mixing theorems in the manuscript. It uses a stopped shell
coupling, square blocks of several sizes, and finite-dimensional linear
algebra. Every probabilistic construction has finite support. No
continuous-time process, infinite-volume measure, or convergence-to-equilibrium
limit is used. A blocking corollary treats every finite interaction range.
\end{abstract}
\maketitle

\section{Scope, conventions, and the exact replacement}
\label{sec:scope}

The conclusion needed in manuscript Theorem 4.10 is a uniform lower bound for
the cell heat-bath gap on sufficiently large square boxes, for both offsets
$o=0,1$. Neither strong mixing on every square nor a gap on arbitrary unions
of compatible squares is used there for any other purpose. We prove the
following replacement. Its window hypothesis is defined precisely in
Section~\ref{sec:shell}; it is exactly manuscript Definition 4.7.

We identify manuscript statements by their source labels in the companion
report. In the supplied PDF the window definition is 4.7 and the surgery
proposition is 4.8; the task's reference to Proposition 4.7 refers to that
same surgery statement, labeled \texttt{p:surgery}.

\begin{theorem}[Window estimate implies uniform box gap]
\label{thm:main}
Let $S$ be a nonempty finite set. Let $\Phi$ be a real-valued,
translation-invariant interaction on $\Z^2$ with
$\Phi_A=0$ if $\operatorname{diam}_\infty A>1$. Suppose its Gibbs
specification satisfies the window estimate with parameters
$l_0\ge1$, $p_0\in(0,1]$, and $M_\star\in\Z_{\ge1}$.
Let $\alpha\in(0,1]$ be a uniform lower bound on every single-site
conditional probability; one explicit choice is given in
Lemma~\ref{lem:spec}. Put
\[
 d=\min\{1/2,p_0\alpha^{16l_0}\}.
\]
Choose any integer $n$ satisfying
\begin{equation}\label{eq:nchoice}
 \begin{split}
 n&\ge44,\qquad n^8\ge M_\star,\\
 n^2&\ge\max\{50000,\,6(l_0+1)\},\\
 n^{12}d^{20}&\ge50000\,2^{20}\,20!.
 \end{split}
\end{equation}
Define
\begin{equation}\label{eq:constants}
 M=n^8,\qquad L=4M+1,\qquad v=L^2,\qquad
 g=\frac{\alpha^{v+1}}{v},\qquad m=\frac{g}{1000}>0.
\end{equation}
For every translate of
$W=\{1,\ldots,N_1\}\times\{1,\ldots,N_2\}$ with $N_1,N_2\ge L$,
every boundary condition $\eta$, and every real function $f$ on $S^W$,
\begin{equation}\label{eq:mainPI}
 \Var_{\mu_W^\eta} f
 \le m^{-1}\sum_{x\in W}\mu_W^\eta(\Var_x f).
\end{equation}
Thus the gap of $\sum_{x\in W}(I-E_x)$ is at least $m$.
There is no parity or divisibility condition on $N_1,N_2$.
\end{theorem}

For example, the integer $n$ can be chosen as the ceiling of the maximum of
\[
44,\quad M_\star^{1/8},\quad \sqrt{50000},\quad
\sqrt{6(l_0+1)},\quad
\bigl(50000\,2^{20}\,20!\,d^{-20}\bigr)^{1/12}.
\]
One may instead take the least integer satisfying the finite list
\eqref{eq:nchoice}. Increasing $n$ preserves each inequality. The constants
are deliberately very conservative.

All gaps in this appendix use \emph{one heat-bath update per site in the
unnormalized Dirichlet form}. Dividing the generator by $|W|$ divides its gap
by $|W|$. The rate-one-per-cell convention in manuscript Section 4 agrees
with ours.

\section{Finite specifications and elementary probability}
\label{sec:finite}

For a finite nonempty $A\subset\Z^2$, an interaction term $\Phi_A$ is a
function $S^A\to\R$. Translation invariance identifies translated terms.
Range $r\in\Z_{\ge1}$ means that all nonzero terms have diameter at most $r$
in the $\ell^\infty$ metric. There are only finitely many possible nonzero
terms containing a given site. Write
\begin{equation}\label{eq:B}
 B=\sum_{A\ni0}\|\Phi_A\|_\infty<\infty,
 \qquad q=|S|,\qquad \alpha_\Phi=q^{-1}e^{-2B}.
\end{equation}
For finite $V$, put
$\partial_r V=\{z\notin V:\dist(z,V)\le r\}$.
For $\eta\in S^{\partial_rV}$ and $\sigma\in S^V$, define
\begin{equation}\label{eq:spec}
 H_V^\eta(\sigma)=\sum_{A:A\cap V\ne\varnothing}
                 \Phi_A((\sigma\eta)_A),\qquad
 \mu_V^\eta(\sigma)=\frac{e^{-H_V^\eta(\sigma)}}
                    {\sum_{\zeta\in S^V}e^{-H_V^\eta(\zeta)}}.
\end{equation}
The sum is finite. The notation $\sigma\eta$ is only used on sites needed by
that sum. In particular, an exterior configuration always means finitely
many letters; writing it on $\Z^2\setminus V$ is harmless shorthand. For
$V=\varnothing$, there is one configuration of probability one.

\begin{lemma}[Consistency, Markov property, and finite energy]
\label{lem:spec}
The probabilities in \eqref{eq:spec} are strictly positive. If $U\subset V$
and $\zeta\in S^{V\setminus U}$, their conditional law on $U$ given
$\sigma_{V\setminus U}=\zeta$ is $\mu_U^{\zeta\eta}$. This conditional law
depends only on the letters in $\partial_rU$.
Every single-site conditional probability is at least $\alpha_\Phi$.
More generally, if all fully specified single-site conditionals have lower
bound $\alpha$, then for any $A\subset V$, any prescribed letters $a\in S^A$,
and any specification of a subset of $V\setminus A$,
\begin{equation}\label{eq:assignment}
 \Pr(\sigma_A=a\mid\hbox{that specification})\ge\alpha^{|A|}.
\end{equation}
All conditioning events here have positive probability.
\end{lemma}
\begin{proof}
Positivity follows from positive real exponentials and a positive finite
normalizing sum. In the conditional law, terms not meeting $U$ are constant
and cancel. The remaining terms are precisely $H_U^{\zeta\eta}$. A term
meeting $U$ cannot see a site more than $r$ away from $U$, proving the Markov
property. Changing a letter at $x$ changes the sum of terms containing $x$
by at most $2B$. Thus, for a proposed letter $s$, every other conditional
weight divided by its weight is at most $e^{2B}$, and its conditional
probability is at least $(qe^{2B})^{-1}$.
A single-site probability conditioned on fewer letters is a finite weighted
average of fully specified conditionals, and hence is at least $\alpha$.
Order the sites of $A$ and apply the ordinary finite chain rule, one letter
at a time, to obtain \eqref{eq:assignment}.
\end{proof}

We will also use product spaces $\Om=\prod_{x\in V}S_x$ with possibly
different nonempty finite alphabets, and a strictly positive probability
$\mu$ on $\Om$. For $U\subset V$ define
\[
 E_Uf(\omega)=\sum_{a\in\prod_{x\in U}S_x}
      \mu(a\mid\omega_{V\setminus U})f(a\omega_{V\setminus U}),
 \qquad \Var_U f=E_U(f^2)-(E_Uf)^2.
\]
Thus $E_U$ resamples the letters of $U$, leaving its complement fixed.
Subscripts $x$ abbreviate $\{x\}$.

\begin{lemma}[Conditional expectation identities]
\label{lem:CE}
With inner product $\langle f,h\rangle_\mu=\mu(fh)$, $E_U$ is a
self-adjoint orthogonal projection; it preserves constants and $\mu$.
Moreover
\begin{equation}\label{eq:projection}
 \langle f,(I-E_U)f\rangle_\mu=\mu(\Var_Uf)
          =\mu\bigl((f-E_Uf)^2\bigr),
\end{equation}
and
\begin{equation}\label{eq:totalvariance}
 \Var_\mu f=\mu(\Var_Uf)+\Var_\mu(E_Uf).
\end{equation}
For a list $\mathscr B=(B_a)_{a\in\mathcal A}$ of subsets of $V$, with
repetitions allowed, put
\[
 H_{\mathscr B}=\sum_{a\in\mathcal A}(I-E_{B_a}),\qquad
 \cD_{\mathscr B}(f)=\sum_a\mu(\Var_{B_a}f).
\]
Then $H_{\mathscr B}$ is self-adjoint and nonnegative, and its quadratic
form is $\cD_{\mathscr B}$. For single sites write $H_V$ and $\cD_V$.
\end{lemma}
\begin{proof}
Group every sum by the fixed value $b$ of the complement of $U$. On that
fiber, $E_U$ is averaging with respect to $\mu(\cdot\mid b)$, so it is
idempotent. Both $\langle f,E_Uh\rangle$ and
$\langle E_Uf,h\rangle$ equal
$\sum_b\mu(b)\mu(f\mid b)\mu(h\mid b)$.
The same grouping proves $\mu E_U=\mu$ and
$\mu(fE_Uf)=\mu((E_Uf)^2)$. Expanding the square proves
\eqref{eq:projection}. Subtract $\mu(f)^2$ from
$\mu(f^2)=\mu(\Var_Uf)+\mu((E_Uf)^2)$ to prove
\eqref{eq:totalvariance}. Sum \eqref{eq:projection} over the list.
\end{proof}

For a nonnegative self-adjoint $H$ annihilating constants, its gap means
\[
 \gap(H)=\inf_{\Var_\mu f>0}
       \frac{\langle f,Hf\rangle_\mu}{\Var_\mu f}.
\]
The infimum over the empty set is $+\infty$. Equivalently a lower bound $c$
means the Poincar\'e inequality
$c\Var_\mu f\le\langle f,Hf\rangle_\mu$ for all $f$.
Bounds for complex functions follow by applying the real bound to their
real and imaginary parts and adding the two inequalities.

For probabilities $p,p'$ on a finite set $X$, define
$\TV(p,p')=\frac12\sum_{a\in X}|p(a)-p'(a)|$.
A coupling is a probability on $X\times X$ with these marginals.

\begin{lemma}[Finite couplings]
\label{lem:coupling}
There is a coupling $(A,A')$ with
$\Pr(A\ne A')=\TV(p,p')$, and every coupling has disagreement probability
at least this value. Marginalization cannot increase total variation.
If $p,p'$ are probabilities on configurations in $V$, and $Q\subset V$
satisfies $\TV(p_Q,p'_Q)\le\varepsilon$, there is a coupling with
\begin{equation}\label{eq:TVHamming}
 \E d_H(\omega,\omega')\le |V\setminus Q|+|V|\varepsilon,
 \qquad d_H(\omega,\omega')=\sum_{x\in V}\1_{\omega_x\ne\omega'_x}.
\end{equation}
\end{lemma}
\begin{proof}
Set $a(t)=\min\{p(t),p'(t)\}$ and
$t_0=\sum_t a(t)=1-\TV(p,p')$.
Put mass $a(t)$ at $(t,t)$. If $t_0<1$, couple the two residual
probabilities independently, with total mass $1-t_0$. Their supports are
disjoint, so the extra mass gives no equality. The cases $t_0=0,1$ are
included by omitting the nonexistent part of the construction.
Conversely the diagonal mass at $(t,t)$ is at most $a(t)$.
For an arbitrary coupling and event $A$, the difference of its two
probabilities is at most the disagreement probability. Alternatively,
summing absolute values over fibers proves marginalization directly.
Couple the $Q$ marginals by the construction just given and complete each
configuration independently using its conditional distribution given $Q$.
On the event of agreement on $Q$ there are at most $|V\setminus Q|$
disagreements, and otherwise at most $|V|$, which gives
\eqref{eq:TVHamming}. Zero-mass conditional fibers are never sampled and
may be assigned any probability.
\end{proof}

\begin{lemma}[Finite surgery]
\label{lem:surgery}
Let $\mu=\mu_V^\tau$, $\mu'=\mu_V^{\tau'}$, $F\subset V$, and let $\nu$
be any coupling of $\mu,\mu'$. For each input pair $(\omega,\omega')$,
choose any coupling of
$\mu_F^{\omega\tau}$ and $\mu_F^{\omega'\tau'}$, replace the letters of
$F$ by this pair, and keep all other letters fixed. The resulting
probability is a coupling of $\mu,\mu'$.
\end{lemma}
\begin{proof}
For the first marginal, summing over the second output gives the conditional
resampling kernel $E_F$, independent of the second input. Summing over
inputs therefore gives $\mu E_F=\mu$, by Lemmas~\ref{lem:spec}
and~\ref{lem:CE}. The second marginal is identical in form. This also
shows why a coupling chosen separately for each input pair is legitimate.
\end{proof}

\section{Positive small-volume gaps and block comparison}
\label{sec:comparison}

\begin{lemma}[An explicit finite-volume bound]
\label{lem:smallgap}
Let $V$ have $t\ge1$ sites, with nonempty finite alphabets $S_x$. Suppose
$\mu>0$ and every single-site conditional probability is at least
$\alpha\in(0,1]$. Then
\begin{equation}\label{eq:smallgap}
 \gap(H_V)\ge\frac{\alpha^{t+1}}{t}.
\end{equation}
Consequently for $1\le t\le v$ the lower bound
$\alpha^{v+1}/v$ holds, uniformly in all boundary conditions satisfying
this same finite-energy bound.
\end{lemma}
\begin{proof}
Lemma~\ref{lem:spec}'s chain-rule argument, which did not use the Gibbs
form, gives $\mu(\sigma)\ge\alpha^t$.
For $s\in S_x$, let $\sigma^{x,s}$ be $\sigma$ with letter $s$ at $x$.
Expanding the conditional variance on each fiber gives
\begin{equation}\label{eq:edgeD}
 \cD_V(f)=\frac12\sum_{\sigma,x,s}
    \mu(\sigma)\mu(s\mid\sigma_{V\setminus x})
          \bigl(f(\sigma)-f(\sigma^{x,s})\bigr)^2.
\end{equation}
Each coefficient before the square is at least $\alpha^{t+1}/2$.
Fix an ordering of $V$. For each ordered pair $(\sigma,\zeta)$ change its
letters to those of $\zeta$ in that order, omitting unchanged letters.
This path has at most $t$ edges. Cauchy--Schwarz gives
\[
 (f(\sigma)-f(\zeta))^2
 \le t\sum_{e\text{ on that path}}(\nabla_e f)^2.
\]
Since
$\Var_\mu f=\frac12\sum_{\sigma,\zeta}
\mu(\sigma)\mu(\zeta)(f(\sigma)-f(\zeta))^2$,
and the total weight of pairs using any particular directed edge is at
most $\sum_{\sigma,\zeta}\mu(\sigma)\mu(\zeta)=1$, we obtain
\[
 \Var_\mu f\le\frac t2\sum_{\text{all directed single-site edges }e}
                (\nabla_e f)^2
 \le t\alpha^{-t-1}\cD_V(f).
\]
This proves \eqref{eq:smallgap}. Since $\alpha\le1$ and $t\le v$, its
right-hand side is at least $\alpha^{v+1}/v$.
\end{proof}

\begin{lemma}[Block comparison, including multiplicities]
\label{lem:comparison}
Let $\mu>0$ on a finite product space, and let $(B_a)_{a\in\mathcal A}$
be a nonempty finite list of nonempty blocks. Suppose that, for every
$a$ and every outside configuration, the single-site heat bath on $B_a$
under that conditional law has gap at least $g>0$. Put
$t_x=\#\{a:x\in B_a\}$ and $t_{\max}=\max_x t_x$.
If $\gap(H_{\mathscr B})\ge c>0$, then
\begin{equation}\label{eq:blockcompare}
 \Var_\mu f\le\frac{t_{\max}}{cg}\cD_V(f).
\end{equation}
\end{lemma}
\begin{proof}
In each fiber outside $B_a$, its Poincar\'e inequality gives
$\Var_{B_a}f\le g^{-1}\sum_{x\in B_a}E_{B_a}(\Var_xf)$.
Here conditioning on the fiber has not changed the one-site conditional
law. Integrate and sum over $a$ to get
\[
 \cD_{\mathscr B}(f)\le g^{-1}\sum_x t_x\mu(\Var_xf)
           \le (t_{\max}/g)\cD_V(f).
\]
Combine this with $c\Var_\mu f\le\cD_{\mathscr B}(f)$.
\end{proof}

\section{The exact window hypothesis and stopped shell coupling}
\label{sec:shell}

Throughout this section the interaction has range one in $\ell^\infty$;
it may have diagonal or many-body terms. Fix
$C_u=\{x\in\Z^2:|x|_\infty\le u\}$, with integer $u\ge1$, an exterior
site $y$ with $\dist(y,C_u)=1$, and set
\[
 s=u^{1/4},\quad k_i=\lfloor is\rfloor,\quad
 i_1=\lfloor s\rfloor+1.
\]
For $i\ge1$ define
\begin{align}\label{eq:shells}
 F_i&=\{x\in C_u:|x-y|_\infty>(i-1)s\},\nonumber\\
 D_i&=\{x\in C_u:is-1<|x-y|_\infty\le is\}
      =\{x\in C_u:|x-y|_\infty=k_i\},\\
 A_i&=\{x\in D_i:\dist(x,\Z^2\setminus C_u)\le l_0\},
 \qquad B_i=D_i\setminus A_i.\nonumber
\end{align}
For $l_0\ge1$ put
\begin{equation}\label{eq:geomscale}
 a(l_0)=\max\{16,6(l_0+1)\},\qquad M_{\rm g}(l_0)=a(l_0)^4.
\end{equation}
Let $I_0$ be the set of indices in $\{2,\ldots,i_1\}$ for which some
corner of $C_u$ is at distance strictly less than $s/3$ from $D_i$.
The distance to the empty set is $+\infty$.

\begin{definition}[Window estimate]
\label{def:window}
The specification has the window estimate with
$(l_0,p_0,M_\star)$ if $l_0\in\Z_{\ge1}$, $p_0\in(0,1]$,
$M_\star\in\Z_{\ge1}$, and for every integer $u\ge M_\star$, every $y$
with $\dist(y,C_u)=1$, every integer $i$ with $2\le i\le s$ and
$i\notin I_0$, and all boundary configurations $\eta,\xi$ for $F_i$,
\begin{equation}\label{eq:window}
 \TV\bigl(\mu_{F_i,B_i}^{\eta},\mu_{F_i,B_i}^{\xi}\bigr)\le1-p_0.
\end{equation}
Here $\mu_{F_i,B_i}$ denotes its marginal on $B_i$.
\end{definition}

For translated squares all definitions are translated with the square.
Translation invariance transfers \eqref{eq:window} to them. The use of the
same $I_0$ as in the manuscript is intentional; no new window shapes are
being requested.

\begin{lemma}[Shell geometry with an explicit threshold]
\label{lem:geometry}
If $u\ge M_{\rm g}(l_0)$, then $I_0$ has at most one element,
$D_i\subset F_i$ for $2\le i\le i_1$, and
\begin{equation}\label{eq:Acount}
 |A_i|\le8l_0\quad(2\le i\le i_1,\ i\notin I_0).
\end{equation}
For $i\ge2$, every site in $(C_u\setminus F_i)$ at distance at most one
from $F_i$ belongs to $D_{i-1}$. The site $y$ is not at distance at most
one from $F_i$.
\end{lemma}
\begin{proof}
We have $s\ge a(l_0)$, $k_i\le s^2+s$, and
\[
 s^2+\tfrac43s<u,\qquad s^2+s<u/2,\qquad
 \lfloor s\rfloor>2s/3,\qquad s/3>l_0.
\]
Each follows immediately from $u=s^4$ and $s\ge16$, except the last,
which also uses $s\ge6(l_0+1)$.
For counting geometry only, interchange axes and reverse their signs so
that $y=(b,-u-1)$, with $-u-1\le b\le u+1$. This does not assume any
reflection symmetry of the interaction. The upper corners have distance
$2u+1$ from $y$. Of the lower corners, at least one has distance at least
$u$ from $y$, because $|b+u|+|b-u|\ge2u$. A corner within $s/3$ of
$D_i$ has distance less than $k_i+s/3\le s^2+4s/3<u$ from $y$.
Thus only one corner can occur for any of these indices. Distinct radii
$k_i,k_j$ differ by at least $\lfloor s\rfloor>2s/3$, so that corner
cannot lie within $s/3$ of two shells. Hence $|I_0|\le1$.

Fix $i\notin I_0$ and $k=k_i$. The bottom horizontal side of the square
$y+[-k,k]^2$ is outside $C_u$. Its top side, if present, is at distance
$k\ge s$ from the bottom exterior half-plane and more than $u$ from the
top one. A vertical side meeting $C_u$ has coordinate $b-k$ or $b+k$.
For the side $b-k$, its distance to the left exterior half-plane is
$e=b-k+u+1\ge1$, while its distance to the right exterior half-plane is
at least $k$. Its lowest point in $C_u$ is $(b-k,-u)$, whose distance
to the lower left corner is $e-1$. Since $i\notin I_0$,
$e-1\ge s/3$. The side $b+k$ satisfies the corresponding right-hand
statement. These facts also cover the endpoints $b=\pm(u+1)$.
Consequently every shell side parallel to an edge of $C_u$ is farther
than $l_0$ from the exterior half-plane beyond that edge. A shell site
within $l_0$ of an edge must therefore lie on a perpendicular shell side.
Each such side contributes at most $l_0$ sites per perpendicular edge.
There are four shell sides and two perpendicular edges per side, proving
\eqref{eq:Acount} by overcounting.

Since $k_i\ge is-1>(i-1)s$, $D_i\subset F_i$.
If $x\in C_u\setminus F_i$ and $z\in F_i$ have $|x-z|_\infty\le1$,
then integrality gives
$|x-y|_\infty=\lfloor(i-1)s\rfloor$, so $x\in D_{i-1}$.
Finally every site of $F_i$ has distance greater than
$(i-1)s\ge s>1$ from $y$.
\end{proof}

\begin{lemma}[Stopped shell coupling]
\label{lem:stopped}
Let the range-one Gibbs specification on the nonempty finite alphabet $S$
have single-site conditional probabilities at least $\alpha\in(0,1]$.
Suppose the window estimate holds with $(l_0,p_0,M_\star)$, and let
\[
 u\ge\max\{M_\star,M_{\rm g}(l_0)\}.
\]
Let $\tau,\tau'$ be boundary conditions for $C_u$ agreeing except
possibly at one exterior site $y$.
Put $\delta=p_0\alpha^{16l_0}$.
If $\dist(y,C_u)>1$, the two measures are equal. Otherwise, for every
integer $h$ with $4\le h\le\lfloor s\rfloor$, there is a coupling of
$\mu_{C_u}^\tau,\mu_{C_u}^{\tau'}$ for which
\begin{equation}\label{eq:stopped}
 \begin{gathered}
 \Pr\bigl(\omega_{Q_h}\ne\omega'_{Q_h}\bigr)
       \le(1-\delta)^{h-3},\\
 Q_h=\{x\in C_u:|x-y|_\infty>hs\}=F_{h+1}.
 \end{gathered}
\end{equation}
In particular the total variation of those marginals has the same bound.
\end{lemma}
\begin{proof}
The first assertion is the Markov property. For $\dist(y,C_u)=1$, start
with the product coupling, and perform the following finite surgeries for
$i=2,\ldots,h+1$. Skip an index in $I_0$.
At a remaining index, if the input pair agrees on $D_{i-1}$, the two
conditional laws on $F_i$ are identical: their interior exterior neighbors
belong to $D_{i-1}$, and the fixed boundaries agree on all other exterior
neighbors, by Lemma~\ref{lem:geometry}. Resample $F_i$ identically in the
two components.

If the pair does not agree on $D_{i-1}$, maximally couple the two
$B_i$ marginals, and complete each sample independently from its
conditional law on $F_i\setminus B_i$. Then independently resample the
letters on $A_i$ in the two components, conditional on the rest of their
respective configurations. Each component still has its required law on
$F_i$: completing a sampled marginal reconstructs its law, and the last
resampling preserves that law by Lemma~\ref{lem:CE}. Thus the operation is
indeed a surgery as in Lemma~\ref{lem:surgery}.
For $i\le h\le s$ its probability of agreement on $B_i$ is at least
$p_0$, by \eqref{eq:window}. Fix any one assignment on $A_i$.
Conditional on everything before the last independent resamplings, both
choose this assignment with probability at least $\alpha^{2|A_i|}$,
by \eqref{eq:assignment}. The probability of agreement on $D_i$ after
step $i$, conditional on any input history, is therefore at least
$p_0\alpha^{2|A_i|}\ge\delta$. This remains true in the identical
resampling branch, where the probability is one. No lower bound is claimed
or needed at step $h+1$ if $h+1>s$.

Let
\[
 J=\{i\in\{2,\ldots,h\}:i\notin I_0,\ i+1\notin I_0\}.
\]
Since at most one exceptional index can remove at most two of the $h-1$
indices, $|J|\ge h-3$. If agreement on $D_i$ occurs after a trial
$i\in J$, step $i+1$ makes the pair identical on $F_{i+1}$.
That agreement persists: a later skipped step changes nothing, and at a
later unskipped step $j$ we have $F_j\subset F_{i+1}$ and
$D_{j-1}\subset F_{i+1}$. Thus its identical-resampling branch is used.
The latter inclusion follows from
$\lfloor(j-1)s\rfloor>is$ when $j\ge i+2$ and $s\ge16$.
Since $Q_h\subset F_{i+1}$, a single successful trial in $J$ ensures final
agreement on $Q_h$.

For completeness, list the trial indices increasingly. Conditional on all
outcomes before any one of them, its failure probability is at most
$1-\delta$. Multiplying these conditional bounds by finite induction
bounds the probability of failure at every trial by
$(1-\delta)^{|J|}\le(1-\delta)^{h-3}$. Lemma~\ref{lem:surgery} ensures
that the final pair has the required marginals. Lemma~\ref{lem:coupling}
then gives the total variation bound. The case $\delta=1$ gives zero
failure probability and is covered by the same induction.
\end{proof}

\begin{corollary}[Manuscript Proposition 4.8, with its original conclusion]
\label{cor:originalprop}
Under the hypotheses of Lemma~\ref{lem:stopped},
\[
 \TV\bigl(\mu_{C_u,Q}^{\tau},\mu_{C_u,Q}^{\tau'}\bigr)
 \le(1-\delta)^{u^{1/4}-4},\qquad
 Q=\{x\in C_u:|x-y|_\infty>u^{1/2}\}.
\]
\end{corollary}
\begin{proof}
Take $h=\lfloor s\rfloor$. Then $hs\le s^2=u^{1/2}$, so $Q\subset Q_h$,
and $h-3\ge s-4$. For $\delta<1$ apply Lemmas~\ref{lem:stopped}
and~\ref{lem:coupling}. If $\delta=1$, the left side is zero and the
right side is zero since $s\ge16$.
\end{proof}

\section{A finite matrix contraction lemma}
\label{sec:contraction}

For a real function on a finite product space let
\[
 \Lip(f)=\max_{d_H(\omega,\omega')=1}|f(\omega)-f(\omega')|,
\]
with value zero if there are no such pairs. Changing disagreeing letters
one at a time proves
$|f(\omega)-f(\omega')|\le\Lip(f)d_H(\omega,\omega')$ for all pairs.
In particular $\Lip(f)=0$ implies that $f$ is constant.

\begin{lemma}[One-step coupling implies an unnormalized gap]
\label{lem:contraction}
Let $\mu>0$ on a finite product space and let
$(B_a)_{a\in\mathcal A}$ be a nonempty finite list of blocks.
Write $A=|\mathcal A|$ and $P=A^{-1}\sum_a E_{B_a}$.
Suppose $0<c\le A$ and that, for every pair of input configurations at
Hamming distance one, their two output laws under $P$ have a coupling
with expected Hamming distance at most $1-c/A$.
Then
\begin{equation}\label{eq:contractgap}
 \gap\left(\sum_a(I-E_{B_a})\right)\ge c.
\end{equation}
\end{lemma}
\begin{proof}
The coupling and the preceding Lipschitz inequality give
$\Lip(Pf)\le(1-c/A)\Lip(f)$ for every real $f$.
By Lemma~\ref{lem:CE}, $P$ is self-adjoint, preserves constants, and is
positive semidefinite. It is also a contraction in $L^2(\mu)$: each
$E_B$ is an orthogonal projection. In finite dimension it has a real
orthonormal eigenbasis. To see the relevant elementary reduction, conjugate
its matrix by $\operatorname{diag}(\sqrt{\mu(\omega)})$; this is a real
symmetric matrix, so orthogonal diagonalization applies. Equivalently one
can minimize its Rayleigh quotient on the finite-dimensional unit sphere
and then induct on its orthogonal complement. No Markov-chain limit is
involved.

For any nonconstant eigenfunction $Pf=\lambda f$ we have $\Lip(f)>0$;
the Lipschitz contraction implies $|\lambda|\le1-c/A$. Hence every
nonconstant eigenvector of $A(I-P)$ has eigenvalue at least $c$. There
cannot be an additional nonconstant eigenvector for eigenvalue one of
$P$. Expand $f-\mu f$ in this finite orthonormal eigenbasis to get
$c\Var_\mu f\le\langle f,A(I-P)f\rangle_\mu$. This is
\eqref{eq:contractgap}. If the product space has one point the Poincar\'e
inequality is automatic.
\end{proof}

\section{Square blocks near the boundary of a rectangle}
\label{sec:count}

The next counting lemma is what allows us to use only square laws, even
at the boundary of the large box. Truncating a translated square at that
boundary would require information on nonsquare regions and is not used.

\begin{lemma}[Clamped squares of varying sizes]
\label{lem:count}
Let $M\ge1$ be an integer, $L=4M+1$, and
$W=\{1,\ldots,N_1\}\times\{1,\ldots,N_2\}$ with $N_1,N_2\ge L$.
Define the finite label set
\[
 \mathcal A=\{M,M+1,\ldots,2M\}
        \times\{1-L,\ldots,N_1\}\times\{1-L,\ldots,N_2\}.
\]
For $a=(u,a_1,a_2)$ set $l=2u+1$ and
\[
 b_j=\min\{\max\{a_j,1\},N_j-l+1\},\qquad
 B_a=\{b_1,\ldots,b_1+l-1\}\times\{b_2,\ldots,b_2+l-1\}.
\]
All $B_a$ are squares contained in $W$, and identical squares with
different labels are retained. For $x\in W$ put
\[
 T_x=\#\{a:x\in B_a\},\qquad
 S_x=\#\{a:x\notin B_a,\ \dist(x,B_a)=1\}.
\]
Then
\begin{equation}\label{eq:counts}
 4M^3\le T_x\le1000M^3,\qquad S_x\le1000M^2.
\end{equation}
The number of labels is
$A=(M+1)(N_1+L)(N_2+L)$.
\end{lemma}
\begin{proof}
Fix a side length $l=2u+1$. The clamping map sends an integer $a_j$ to
the interval of allowable starting points $[1,N_j-l+1]$.
For each of the $l$ integers
$a_j\in[x_j-l+1,x_j]$, the clamped interval of length $l$ still contains
$x_j$. Indeed an upward clamp starts at $1\le x_j$ and only moves the
right endpoint to the right; a downward clamp ends at $N_j\ge x_j$ and
only moves the left endpoint to the left. These integers all belong to
the label range. Therefore $T_x\ge\sum_{u=M}^{2M}(2u+1)^2
\ge4M^3$.

A useful upper bound is that an interval of $t$ allowable starting points
has at most $t+2L$ preimages under clamping. Away from the endpoints there
is one preimage per point; the left endpoint has at most $L$ extra
preimages and the right endpoint at most $l-1\le L$ extra preimages.
This bound remains valid if the two endpoints coincide.
The starting points whose interval contains $x_j$ form an interval with
at most $l$ points; a slightly sharper count is
$l+L+(l-1)\le3L$. Thus
\[
 T_x\le(M+1)(3L)^2\le450M^3\le1000M^3,
\]
using $M+1\le2M$ and $L\le5M$.

For the bound on $S_x$, a square at exterior distance one from $x$ has,
in at least one coordinate $j$, a starting point equal to
$x_j+1$ or $x_j-l$. Call this its directed side, overcounting corners.
In the other coordinate the interval must be within distance one of
$x$, so its allowable starting points form an interval of length at most
$l+2$. Its number of nominal labels is at most
$2l+L+1\le4L$.

If the prescribed starting point in coordinate $j$ is strictly between
$1$ and $N_j-l+1$, it has one preimage. These regular cases contribute
at most $4(M+1)4L=16(M+1)L$ labels.
The equality $x_j+1=1$ is impossible for $x\in W$, and
$x_j-l=N_j-l+1$ would give $x_j=N_j+1$, also impossible.
The remaining endpoint cases require respectively
\[
 l=N_j-x_j\quad\hbox{or}\quad l=x_j-1.
\]
For each of the four directed sides, there is at most one $u$ with
that side length. The constrained coordinate has at most $2L$ preimages
at an endpoint, including the case of coincident endpoints. Together
with the bound $4L$ in the other coordinate, these at most four
exceptional cases contribute at most $32L^2$ labels. We have proved
\[
 S_x\le16(M+1)L+32L^2\le160M^2+800M^2\le1000M^2.
\]
The cardinality of $\mathcal A$ follows by multiplying its three lengths.
\end{proof}

\begin{definition}[A finite family of boundary-coupling bounds]
\label{def:Kcriterion}
For an integer $M\ge1$, say that $\mathcal K(M,K)$ holds if, for every
integer $u\in[M,2M]$, every translate $B$ of $C_u$, every site
$y\notin B$ with $\dist(y,B)=1$, and every pair of boundary conditions
agreeing off $y$, the Gibbs measures on $B$ have a coupling with
expected Hamming distance at most $K$.
\end{definition}

\begin{theorem}[A direct finite-square criterion]
\label{thm:criterion}
Let a strictly positive range-one specification on $\Z^2$ have uniform
single-site conditional lower bound $\alpha$. Translation invariance is
not needed if the following hypothesis is imposed on every translate.
Suppose
\begin{equation}\label{eq:Ksmall}
 \mathcal K(M,K)\quad\hbox{holds with}\quad K\le M/1000.
\end{equation}
Put $L=4M+1$, $v=L^2$, and $g=\alpha^{v+1}/v$.
For every rectangular box with both sides at least $L$, every boundary
condition, and every $f$,
\begin{equation}\label{eq:criterionPI}
 \Var_\mu f\le(1000/g)\cD_W(f).
\end{equation}
\end{theorem}
\begin{proof}
Use the square blocks of Lemma~\ref{lem:count} and choose their labels
uniformly to define $P=A^{-1}\sum_a E_{B_a}$. Consider two input
configurations differing at one site $x\in W$. If $x\in B_a$, their
conditional block laws are identical; use the same resampled configuration
in both. The updated pair then agrees everywhere. If $x\notin B_a$ and
$\dist(x,B_a)>1$, the Markov property again gives identical block laws;
use the same sample, leaving exactly the original disagreement. In the
remaining $S_x$ cases, the two block laws have boundary conditions
differing only at $x$. Use \eqref{eq:Ksmall}'s coupling. The outside
contributes the original disagreement and the expected inside
contribution is at most $K$.

Using the same uniformly chosen label in both components gives a coupling
of the two $P$ output laws with
\[
 \E d_H\le1-\frac{T_x-KS_x}{A}
 \le1-\frac{M^3}{A},
\]
because $T_x-KS_x\ge4M^3-M^3=3M^3\ge M^3$.
Here $M^3\le A$ follows directly from the formula for $A$.
Lemma~\ref{lem:contraction} gives a block gap at least $M^3$.
Every conditional law on every $B_a$ has at most $v$ sites and the same
finite-energy bound; its single-site gap is at least $g$ by
Lemma~\ref{lem:smallgap}. Finally $t_{\max}\le1000M^3$ in
Lemma~\ref{lem:comparison}, proving \eqref{eq:criterionPI}.
All boundary configurations of the large box were arbitrary throughout.
\end{proof}

\section{Closing the constants and proving the replacement theorem}
\label{sec:scale}

\begin{lemma}[A purely finite tail bound]
\label{lem:tail}
For $0<d\le1/2$ and integer $n\ge44$,
\begin{equation}\label{eq:tail}
 (1-d)^{n-3}\le\frac{2^{20}\,20!}{d^{20}n^{20}}.
\end{equation}
\end{lemma}
\begin{proof}
Set $N=n-3\ge41$. Since $(1-d)^{-1}\ge1+d$, the finite binomial
expansion gives
\[
 (1-d)^{-N}\ge(1+d)^N\ge\binom N{20}d^{20}
 \ge\frac{(N-19)^{20}}{20!}d^{20}
 \ge\frac{(n/2)^{20}}{20!}d^{20}.
\]
All quantities are positive, so taking reciprocals proves the claim.
\end{proof}

\begin{lemma}[The window estimate supplies the square criterion]
\label{lem:windowK}
Under the hypotheses and choices of Theorem~\ref{thm:main},
$\mathcal K(M,M/1000)$ holds for $M=n^8$.
\end{lemma}
\begin{proof}
Fix $u\in[M,2M]$. Then
$s=u^{1/4}\ge n^2\ge6(l_0+1)$ and $s\ge16$; hence
$u\ge M_{\rm g}(l_0)$ and $u\ge M_\star$.
Choose $h=n$ in Lemma~\ref{lem:stopped}, which is allowed since
$4\le n\le\lfloor s\rfloor$. The radius of the exceptional near region
is $ns\le2n^3$, since $u\le2n^8$. Its number of sites is at most
$(4n^3+1)^2\le25n^6$. The square itself has at most
$(4n^8+1)^2\le25n^{16}$ sites. Lemma~\ref{lem:coupling} (or the actual
coupling in Lemma~\ref{lem:stopped}) therefore gives expected Hamming
cost at most
\[
 K_n=25n^6+25n^{16}(1-d)^{n-3}.
\]
The second line of \eqref{eq:nchoice} gives
$25n^6\le n^8/2000$.
Lemma~\ref{lem:tail} and the final inequality of \eqref{eq:nchoice}
give
\[
 25n^{16}(1-d)^{n-3}
 \le\frac{25\,2^{20}\,20!}{d^{20}n^4}
 \le n^8/2000.
\]
Consequently $K_n\le n^8/1000=M/1000$.
Since $d\le\delta$, replacing $\delta$ by $d$ only enlarged the bound.
Translation invariance transfers this conclusion to every translated
square, as required in Definition~\ref{def:Kcriterion}.
\end{proof}

\begin{proof}[Proof of Theorem~\ref{thm:main}]
Lemma~\ref{lem:spec} supplies consistency, the range-one Markov property,
and the required uniform finite-energy bound. The choice
\eqref{eq:nchoice} exists by its displayed explicit formula.
Lemma~\ref{lem:windowK} proves the hypothesis of
Theorem~\ref{thm:criterion}; its conclusion is exactly
\eqref{eq:mainPI} with the constants \eqref{eq:constants}.
\end{proof}

\begin{remark}[Why stopping sooner matters]
\label{rem:stopping}
The manuscript's Proposition 4.8 stops at distance $u^{1/2}$. Bounding
Hamming cost from that statement alone allows order $u$ disagreements,
which is too large for the elementary contraction above. The actual window
hypothesis allows stopping after $n$ shells when $u\asymp n^8$.
The near region then has order $n^6=o(n^8)$ sites, while its failure
probability is bounded by a finite binomial calculation. This is a
strengthening of the coupling proof using exactly the same input, rather
than a deduction of a new mixing assumption from Proposition 4.8's
statement alone.
\end{remark}

\section{Arbitrary finite interaction range}
\label{sec:range}

\begin{corollary}[Blocking a finite-range model]
\label{cor:range}
Let $S$ be nonempty finite, and let $\Phi$ be a real,
translation-invariant interaction of integer range $r\ge1$ on $\Z^2$.
Let $\alpha=q^{-1}e^{-2B}$ be as in \eqref{eq:B}.
Partition the lattice into cells
\[
 C_c=rc+\{0,\ldots,r-1\}^2,\qquad c\in\Z^2,
\]
and regard each cell as one letter of $\widehat S=S^{\{0,\ldots,r-1\}^2}$.
Suppose this cell specification satisfies Definition~\ref{def:window}
with parameters $(l_0,p_0,M_\star)$. Define
$\widehat\alpha=\alpha^{r^2}$ and use
Theorem~\ref{thm:main}'s formulas with $\widehat\alpha$ to obtain
$\widehat L$ and $\widehat m$.
Then every translate of a rectangular box whose side lengths are
$rN_1,rN_2$, with $N_1,N_2\ge\widehat L$, and every boundary condition
satisfy
\begin{equation}\label{eq:rangegap}
 \gap(H_{\rm original})\ge
       \frac{\alpha^{r^2+1}}{r^2}\widehat m>0.
\end{equation}
In particular this covers arbitrarily large squares of side divisible by
$r$, which is sufficient whenever one is free to choose the box side.
\end{corollary}
\begin{proof}
A nonzero interaction can meet only cells whose indices have pairwise
$\ell^\infty$ distance at most one. Indeed, if two cell indices differed
by at least two in one coordinate, all sites in these two cells would be
separated by at least $r+1$ in that coordinate. Regrouping original
interaction terms by the set of cells they meet therefore gives a
translation-invariant, real, range-one interaction on the finite alphabet
$\widehat S$. There are finitely many terms meeting a cell, and the
finite-volume Gibbs probabilities agree exactly under this identification.
By the chain-rule argument \eqref{eq:assignment}, every conditional cell
letter has probability at least $\widehat\alpha$.
Theorem~\ref{thm:main} gives a cell heat-bath gap at least
$\widehat m$ on the indicated rectangles. In any fixed outside fiber, the
original single-site dynamics inside one cell has $r^2$ sites and all
conditionals at least $\alpha$; Lemma~\ref{lem:smallgap} gives gap at
least $\alpha^{r^2+1}/r^2$. Apply
Lemma~\ref{lem:comparison} to the disjoint cells: the multiplicity is one.
This proves \eqref{eq:rangegap}. Translation invariance identifies any
translated cell partition with the original one.
\end{proof}

This corollary imposes the window condition on the actual range-one blocked
specification. Finite energy and finite range alone do not imply a
uniform gap, and no such unconditional claim is being made. For the
manuscript the specification has already been blocked into its $2\times2$
physical cells and already has range one, so there is no additional window
hypothesis or further blocking in the application.

\section{Insertion into manuscript Section 4}
\label{sec:application}

\begin{corollary}[Replacement for manuscript Corollary 4.9]
\label{cor:twooffsets}
Assume exactly the manuscript's window estimate for both cell models
$\widehat\mu^{(0)}$ and $\widehat\mu^{(1)}$, with their finite alphabets
of cardinality $57600$ and real range-one potentials.
Then there are explicitly determined integers $L_0^{(o)}\ge1$ and
numbers $m_0^{(o)}>0$ such that every translated square of $t\times t$
cells, with $t\ge L_0^{(o)}$, has cell heat-bath gap at least $m_0^{(o)}$,
under every boundary condition. The same holds for rectangles with both
sides at least $L_0^{(o)}$.
Set $L_{\rm cell}=\max_o L_0^{(o)}$ and
$m_0=\min_o m_0^{(o)}$. These constants work simultaneously for both
offsets.
\end{corollary}
\begin{proof}
For each of the two finite potentials calculate $B_o$ as in \eqref{eq:B}
and take $\alpha_o=57600^{-1}e^{-2B_o}$. Use that model's given window
parameters in \eqref{eq:nchoice} and \eqref{eq:constants}, setting
$L_0^{(o)}=4(n_o)^8+1$ and
$m_0^{(o)}=\alpha_o^{(L_0^{(o)})^2+1}/
[1000(L_0^{(o)})^2]$. Apply Theorem~\ref{thm:main} twice.
The stated maximum and minimum are over two values and retain the
claimed properties.
\end{proof}

\begin{lemma}[Comparison with physical-site heat bath in the manuscript]
\label{lem:physical}
For the two checkerboard frame models, let $\beta>0$ be the minimum of
all one-physical-site conditional letter probabilities, over the two site
types, all their possible neighbor letters, and all letters at the site.
This is a positive finite minimum. Put
\[
 c_0=\beta^5/4>0.
\]
On every finite union of physical $2\times2$ cells, under every boundary
condition, the physical single-site heat-bath gap is at least $c_0$ times
the cell heat-bath gap.
\end{lemma}
\begin{proof}
Every site weight and every allowed bond weight of the frame model is
strictly positive (see the manuscript equation labeled
\texttt{e:icodod} in the source). There are finitely many choices at the
site and its four neighbors. Thus every conditional probability is
positive, and taking the finite minimum defining $\beta$ proves
$\beta>0$. Equivalently, if $w_-,w_+$ are the smallest and largest site
weights, $b_-,b_+>0$ the smallest and largest allowed bond weights, and
$q_{\max}=20$, one may take the explicit lower bound
\[
 \beta\ge\frac{w_-b_-^4}{20w_+b_+^4}>0.
\]
Condition on everything outside one $2\times2$ cell. Its four-site chain
satisfies Lemma~\ref{lem:smallgap}, so its gap is at least $\beta^5/4$.
Apply Lemma~\ref{lem:comparison} to the disjoint cells, with multiplicity
one. This is also a full finite-state proof of manuscript Lemma 4.2's
comparison.
\end{proof}

\begin{corollary}[The box input used by manuscript Theorem 4.10]
\label{cor:boxinput}
Use the constants $L_{\rm cell},m_0,c_0$ above. Choose an integer
\begin{equation}\label{eq:chooseell}
 t\ge\max\{L_{\rm cell},\lceil1/(c_0m_0)\rceil,1\},
 \qquad \ell=2t.
\end{equation}
For all even torus lengths $L,N\ge\ell+2$, every translated
$\ell\times\ell$ physical box, and every configuration outside it,
the conditional physical single-site heat-bath gap is at least
$c_0m_0$. In the notation of manuscript Section 3,
\begin{equation}\label{eq:gell}
 g(\ell)\ge c_0m_0,\qquad \ell\ge2/(c_0m_0).
\end{equation}
Consequently the manuscript's finite-size criterion gives
\[
 \lambda(\mu_D)\ge c_0m_0/2,
\]
and its already separate frame-to-quantum comparison gives its original
Theorem 4.10 conclusion with
$c=g_\star c_0m_0/4$ and torus threshold $\ell+2$.
\end{corollary}
\begin{proof}
Lift a translated box to $\Z^2$. Because $L,N$ are even and at least
$\ell+2$, its sites are distinct on the torus and no wrap-around bond
joins two lifted box sites. Assign every exterior nearest-neighbor site
of the lifted box the letter of its torus image. This is consistent even
if images repeat. The conditional product of site and bond weights is
then precisely the plane box Gibbs law with these boundary letters;
all other exterior letters can be arbitrary.

The lower left corner of the box has either equal or unequal coordinate
parities. If its second coordinate is odd, translating by $(1,1)$ makes
it even and preserves the checkerboard pattern. Next translate by even
amounts in each coordinate to identify its cells with the offset $o=0$
if the two original parities were equal, and $o=1$ if they were unequal.
The box is now exactly a $t\times t$ square of the corresponding cells.
There are no parity restrictions on $t$.
Corollary~\ref{cor:twooffsets} and Lemma~\ref{lem:physical} give gap
$c_0m_0$, uniformly over its boundary condition. Taking the minimum over
placements and outside configurations proves \eqref{eq:gell}.
The two final displayed gap conclusions are substitutions into the
manuscript's Theorem 3.1 (label \texttt{t:K}) and its frame comparison;
no quantum statement is used to prove any lemma of this appendix.
\end{proof}

\begin{remark}[What is replaced and what has been proved]
Theorem~\ref{thm:main} replaces the \emph{combination} of manuscript
Theorems 4.3--4.4 and their application in Corollary 4.9. We do not assert
the full original Theorem 4.3 conclusion from its one-scale hypothesis,
or the full arbitrary-compatible-set conclusion of Theorem 4.4.
Those stronger assertions may be deleted from the proof's dependency
chain. Corollary~\ref{cor:twooffsets} is their exact box-gap replacement,
and Corollary~\ref{cor:boxinput} shows that the conclusion of Theorem 4.10
is unchanged. In its proof one deletes the integer
$K=(2M_0^{(0)}+1)(2M_0^{(1)}+1)$ and the compatible-set divisibility
argument, and chooses $t,\ell$ by \eqref{eq:chooseell} instead.
Remark 4.5's informal extension from two spin values to $57600$ letters
is no longer required. All arguments here work on an arbitrary finite
alphabet and use no order or attractivity.
\end{remark}

\section{Sources and independence of the proof}
\label{sec:sources}

The shell construction is adapted from \cite[Section 2, proof of
Proposition 2.1, equations (2.4)--(2.12)]{MOS}; its stopped version is
proved here because the shorter stopping radius is essential.
Our window hypothesis is the manuscript's Definition 4.7; its separate
Section 5 proves it from that manuscript's Hypothesis (P). This appendix
proves the implication from that input and introduces no extra model
estimate.
The square-block construction and its elementary counting proof are
provided here in Lemma~\ref{lem:count}; they are not invoked as a result
of another paper.

The comparison in Lemma~\ref{lem:comparison} is the finite conditional
variance argument also found in \cite[Section 3.7, Proposition 3.4,
(3.29)--(3.32)]{Martinelli}. That source's Section 4.2, Theorem 4.5,
gives a particularly useful recursive route under mixing on fat
rectangles. Its hypotheses are on more region shapes than the square
window input in this task, so the direct construction here avoids having
to transfer that hypothesis.
The general coupling viewpoint is also developed in
\cite[Theorem 2.5 and Section 6]{DSVW}.
No rapid-mixing theorem or comparison theorem from that work is used as a
black box.

Other routes considered were Cesi's entropy quasi-factorization
\cite[Section 2, Proposition 2.1, and Section 4, Theorem 4.1]{Cesi} and
the recursive variance proof of
\cite[Section 3.3, Lemma 3.1 and proof of Theorem 2.2]{BCC}. They are useful when the
appropriate ratio or conditional mixing estimates are already known.
They are not dependencies here. Likewise the percolation construction of
\cite[Section 3, Theorem 3.1]{Ott} is a later proof of the full weak-to-strong
mixing implication, rather than a shorter way to the specific finite-box
gap needed here.

The original paper \cite[Section 4, (4.1)--(4.5), Theorem 4.1]{MO}
explains both the heat-bath convention and the danger of truncating square
updates to other shapes. Its general-system argument uses couplings and
recursive time bounds. The finite positive matrix $P$ in
Lemma~\ref{lem:contraction} replaces the time-evolution part entirely.
No statement from \cite{MO} or \cite{MOS} is assumed in the proof.

\begin{thebibliography}{9}
\bibitem{MOS}
F. Martinelli, E. Olivieri, and R. H. Schonmann,
\emph{For 2-D lattice spin systems weak mixing implies strong mixing},
Comm. Math. Phys. \textbf{165} (1994), 33--47.
\href{https://doi.org/10.1007/BF02099735}{doi:10.1007/BF02099735}.
Project Euclid: \url{https://projecteuclid.org/euclid.cmp/1104271032}.
Full author version: \url{https://web.ma.utexas.edu/mp_arc-bin/mpa?yn=93-154}.

\bibitem{MO}
F. Martinelli and E. Olivieri,
\emph{Approach to equilibrium of Glauber dynamics in the one phase region.
I. The attractive case}, Comm. Math. Phys. \textbf{161} (1994), 447--486.
\href{https://doi.org/10.1007/BF02101929}{doi:10.1007/BF02101929}.
Project Euclid: \url{https://projecteuclid.org/euclid.cmp/1104270006}.
Full author version: \url{https://web.ma.utexas.edu/mp_arc-bin/mpa?yn=92-186}.

\bibitem{Martinelli}
F. Martinelli, \emph{Lectures on Glauber dynamics for discrete spin models},
in Lectures on Probability Theory and Statistics, Saint-Flour 1997,
Lecture Notes in Math. \textbf{1717}, Springer, 1999, 93--191.
\href{https://doi.org/10.1007/978-3-540-48115-7_2}{doi:10.1007/978-3-540-48115-7\_2}.
Full lecture preprint:
\url{https://web.ma.utexas.edu/mp_arc-bin/mpa?yn=97-578}.

\bibitem{DSVW}
M. Dyer, A. Sinclair, E. Vigoda, and D. Weitz,
\emph{Mixing in time and space for lattice spin systems: a combinatorial
view}, Random Structures Algorithms \textbf{24} (2004), 461--479.
\href{https://doi.org/10.1002/rsa.20004}{doi:10.1002/rsa.20004}.
Author version: \url{https://people.eecs.berkeley.edu/~sinclair/st_full.pdf}.

\bibitem{Cesi}
F. Cesi, \emph{Quasi-factorization of the entropy and logarithmic Sobolev
inequalities for Gibbs random fields}, Probab. Theory Related Fields
\textbf{120} (2001), 569--584.
\href{https://doi.org/10.1007/PL00008792}{doi:10.1007/PL00008792}.
Author version: \url{https://web.ma.utexas.edu/mp_arc-bin/mpa?yn=00-284}.

\bibitem{BCC}
L. Bertini, N. Cancrini, and F. Cesi,
\emph{The spectral gap for a Glauber-type dynamics in a continuous gas},
Ann. Inst. H. Poincar\'e Probab. Statist. \textbf{38} (2002), 91--108.
\url{https://www.numdam.org/item/AIHPB_2002__38_1_91_0/}.

\bibitem{Ott}
S. Ott, \emph{A new perspective on the equivalence between Weak and Strong
Spatial Mixing in two dimensions}, arXiv:2502.04117.
\url{https://arxiv.org/abs/2502.04117}.
\end{thebibliography}
\end{document}
